This section answers RQ3: \emph{absent a comprehensive benchmark, can lightweight evaluation heuristics and a small set of recurring optimization patterns reliably flag performance-poor DSL kernels and guide their repair?}
RQ1 (\cref{sec:evaluation}) showed that a comprehensive performance benchmark is infeasible and that correctness gates admit slow kernels; RQ2 (\cref{sec:analysis}) explained why the gaps arise.
We now distill that evidence into developer-facing guidance: a two-part \emph{evaluation heuristic} for deciding whether a kernel is efficient (\cref{sec:mitig:heuristics}), and the \emph{optimization patterns} that repair the gaps it flags (\cref{sec:mitig:norm} onward), validated by optimizing kernels across the suite.
The optimization campaigns (\cref{sec:meth:optimization}) ran on a separate development GPU; their kernels are re-timed here on the A100-SXM4 and GH200.
The central result is a clean dichotomy: most of the dramatic gaps are \emph{authoring artifacts} that one dominant pattern fully repairs, while a minority are \emph{genuine residuals} that survive best-effort optimization---and the heuristic tells the two apart without a ground-truth benchmark.

\subsection{Evaluation Heuristics: Is This Kernel Efficient?}
\label{sec:mitig:heuristics}

We propose two complementary checks a developer can apply without a curated benchmark.

\paragraph{A comparability screen (pragmatic, baseline-relative)}
An efficient DSL kernel should (i) come within a small factor of the vendor/PyTorch baseline on representative shapes, (ii) be at least at parity---ideally faster---at large input sizes where launch and framework overheads amortize, and (iii) show no catastrophic per-shape collapse.
This screen is cheap and catches gross failures: the idiomatic TileLang LayerNorm that runs more than $300\times$ slower (\cref{sec:eval:passesslow}) fails it immediately.
Its limitation is circularity---PyTorch is simultaneously the baseline and the de-facto definition of ``good,'' so the screen cannot certify a kernel that merely matches an already-slow baseline, and it can be fooled into crediting a kernel that beats an \emph{unfused} eager path for the wrong reason (the RMSNorm $873\%$ ``win'' in \cref{tab:summary} is a fusion artifact, not headroom).
It is therefore a screen, not a certificate.

\paragraph{A roofline anchor (baseline-independent)}
To break that circularity we anchor quality to the hardware rather than to PyTorch: for a kernel whose essential work is $W$ (bytes moved if memory-bound, FLOPs if compute-bound) and whose optimized time is $t$, the achieved roofline fraction is $\rho = W / (t \cdot P)$, where $P$ is the relevant hardware peak~\cite{williams2009roofline}.
TritonBench reports an analogous achieved-peak metric~\cite{li2025tritonbench}; we use it not as a leaderboard score but as a \emph{decision signal}---a kernel near the memory or compute roof is efficient \emph{regardless} of what the library does.
\Cref{tab:roofline} reports $\rho$ for the optimized kernels on the A100-SXM4 (primary) and the GH200, alongside their PyTorch-relative efficiency $E_\text{lib}$ and the \emph{cliff} (naive-to-optimized speedup) that measures how much authoring headroom each kernel hid; the two architectures agree on every authoring-artifact vs.\ structural-residual classification (the GH200 $\rho$ values cited below are matched within $\pm0.1$ on the A100).
The anchor does two things the comparability screen cannot.
First, it confirms the genuine residuals are genuinely far from the hardware: optimized Triton Conv2d reaches only $\rho=0.16$ and index-returning argmax only $\rho=0.10$, so their shortfall is real headroom, not merely a fast baseline.
Second, it de-inflates baseline-relative outliers: the RMSNorm kernel that looks like a $13\times$ ``win'' over PyTorch sits at $\rho=0.69$---efficient, but not $13\times$ beyond what the hardware allows.
On the GH200, the recovered normalization and reduction family lands at $\rho=0.41$--$0.87$, while the residual convolution and argmax kernels sit at $\rho\le0.28$ even after best-effort optimization (the A100 mirrors the split---$\rho=0.42$--$0.89$ recovered, $\le0.34$ residual---so the classification is architecture-independent); the cliff column shows these are not the same axis---LayerNorm hid a $1541\times$ authoring artifact that the dominant pattern fully recovers, whereas Conv2d's $8.2\times$ cliff still leaves it at $\rho=0.28$, a structural ceiling the rewrite cannot lift.
We present the two checks together because neither suffices alone, and they disagree exactly in a judgment band near $\rho\approx0.5$: Softmax and Matmul both achieve $\rho=0.47$, yet Softmax is at parity ($E_\text{lib}=0.87$) while Matmul trails cuBLAS ($E_\text{lib}=0.68$)---only the two checks together make the call.

% [tab:roofline relocated to appendix (App. Threats, by the Heuristic-assumptions paragraph) to fit the 10-page limit]

\subsection{Normalization Kernels: Correcting RC0}
\label{sec:mitig:norm}

\medskip\noindent\textbf{Finding: Two user-space fixes---\texttt{T.serial}$\to$\texttt{T.reduce} (dominant) and native bf16/fp16 I/O without intermediate \texttt{.float()} casts---bring TileLang LayerNorm to $124\%$ and RMSNorm to $990\%$ of PyTorch on the A100-SXM4 ($1347\times$ and $1157\times$ over the unoptimized baseline), confirming RC0 as the primary, fully user-correctable cause of the normalization deficit.}

The dominant fix is Step~1: a single \texttt{T.reduce} replacing the manual \texttt{T.serial} loop (subsequent iterations confirmed the loop structure, not thread configuration, was the bottleneck).
Step~2, native bfloat16 I/O removing intermediate \texttt{.float()} casts, then brings LayerNorm from 364~ms to 0.270~ms (124\% of PyTorch) and RMSNorm to 0.254~ms (990\%, exceeding 100\% because the fixed kernel fuses the scaling pass that PyTorch's eager two-pass path does not)---near the $\approx$0.18~ms memory-bandwidth floor for an $8192^2$ bfloat16 tensor at the A100's $\approx$1.5~TB/s.
This is not a new technique (\texttt{T.reduce} already existed); the contribution is showing RC0 fully explains the anomaly and is mechanically correctable in user space without algorithmic change.
The full 18-iteration LayerNorm search trajectory is given in \cref{tab:layernorm:trajectory}.

\subsection{Convolution: Partial Recovery via Implicit GEMM}
\label{sec:mitig:conv}

\medskip\noindent\textbf{Finding: Restructuring Conv2d as an FP16 Tensor-Core implicit GEMM with 16-byte-aligned input padding (RC1) and an expanded autotune space (RC2) reaches 36--77\% of cuDNN across the $1\times1$--$7\times7$ filter family on the A100-SXM4 (all correct)---narrowing but not closing the gap.}

The restructuring unfolds the convolution into an on-the-fly implicit GEMM~\cite{zhou2021implicit} that FP16 Tensor Cores accelerate, with 16-byte input padding enabling \texttt{LDG.128} loads (RC1) and an expanded autotune space (smaller \texttt{BLOCK\_K}, more pipeline stages) covering configurations absent from the default list (RC2).
Of the residual $\approx$20\% gap, absent Winograd selection contributes only $\approx$2--3\% (a 0.03\% deterministic-vs-non-deterministic delta for $3\times3$ stride-1); the remainder reflects general cuDNN implicit-GEMM tuning---kernel selection, split-K, shared-memory tile sizing---that the single-lowering DSL kernel does not match (RC4).
Per-filter optimized efficiencies range from 77.4\% at $1\times1$ to 36.1\% at $7\times7$ (all correct); closing this tuning gap and adding in-DSL Winograd support are future work.

\subsection{GEMM and Reduction Kernels}
\label{sec:mitig:other}

Beyond the normalization kernels, we optimized the GEMM kernel and the full TileLang
reduction/softmax family to test how broadly---and how completely---the RQ1 gaps recover.

\textbf{Square matmul (Triton, FP16).}
Adding \texttt{@triton.autotune} with an expanded set of tile configurations
and a \texttt{GROUP\_SIZE\_M} L2 cache swizzle (which reorders output tiles to improve L2 data reuse across the $M$-dimension)
reduces latency from 1.08~ms to 0.79~ms ($1.37\times$, $E_\text{lib} = 77\%$)
on a $4096\times4096$ matmul, re-timed on the A100-SXM4;
the same expanded search recovers $E_\text{lib}$ from $59.7\%$ to $81.9\%$ ($1.37\times$) at the RQ1 $16384^2$ shape.
On the data-center A100, cuBLAS is fast enough that this gap narrows but does not close.
Iterations 2--6 (persistent kernel, expanded configs, \texttt{max\_num\_imprecise\_acc}) converged without further gain,
confirming that the expanded configuration set and L2 swizzle are the effective ceiling for this shape.
This confirms RC2 as a contributor to the GEMM gap
and demonstrates that search-space expansion substantially narrows it for square shapes.
This reconciles with the earlier default-grid finding that the stock 12-config autotune grid recovers $\Delta\approx0$pp:
the gain here is a property of the \emph{expanded} search space
(sweeping \texttt{GROUP\_SIZE\_M}, \texttt{num\_warps}, and \texttt{num\_stages}), not of a different problem shape.

\textbf{Argmax (TileLang, FP16).}
Replacing global-memory element accesses with tiled shared-memory bulk loads
and native FP16 I/O reduces latency from 11.97~ms to 2.31~ms ($5.2\times$) on a $8192\times32768$
reduction, re-timed on the A100-SXM4.
The optimized kernel reaches $E_\text{lib} = 27\%$: it removes most of the unoptimized-TileLang overhead,
but the A100's native \texttt{torch.argmax} (0.62~ms) remains $3.7\times$ faster, so parity is not reached on this GPU.
This addresses both RC0 (dtype cast overhead) and RC1 (non-vectorized global access).

\textbf{Reduction and softmax family (TileLang).}
The same RC0 fix---replacing \texttt{T.serial} reduction loops with native \texttt{T.reduce}
and streaming the reduction dimension in tiles (an online, max-rescaled pass for the softmax
variants) rather than materializing the full row in a fragment---generalizes across the
\emph{entire} TileLang reduction and normalization family on the A100-SXM4 (\cref{tab:mitigation}), not
just a handful of kernels. Every catastrophically slow kernel recovers to
PyTorch-comparable or better: \texttt{max\_reduction}, \texttt{mean\_reduction}, and
\texttt{logsumexp} \emph{exceed} PyTorch ($132\%$, $99\%$, $582\%$), and \texttt{softmax} and
\texttt{log\_softmax} reach $86\%$ and $71\%$---all numerically correct, with $4$--$29\times$
speedups over the idiomatic kernels. The lone within-A100 exception is
index-returning \texttt{argmax} ($27\%$), where the index bookkeeping is intrinsically heavier
and \texttt{torch.argmax} is already well tuned---note that the \emph{value-only} reduction over
the identical shape (\texttt{max\_reduction}) reaches $132\%$, isolating the residual to the index pass.
Re-timing the same optimized kernels on the GH200 confirms the pattern transfers---LayerNorm $120\%$,
RMSNorm $1303\%$, MeanReduction $97\%$, MaxReduction $84\%$, Softmax $87\%$, LogSoftmax $90\%$
(\cref{tab:mitigation,tab:roofline})---with one instructive exception: \texttt{logsumexp}, whose
A100-tuned wide fp32 fragment register-spills on \texttt{sm\_90} (255~registers, ${\sim}280$~GB of
local-memory spill traffic, $12\%$ occupancy by Nsight Compute), collapsing to $1.0\%$. The spill is an RC3-class effect, not
the RC0 reduction idiom: retuning the fragment width recovers it only to $63\%$, so on the GH200
\texttt{logsumexp} is a genuine residual---an ``optimized'' kernel validated on one GPU that is
$100\times$ off on another with no correctness signal, precisely the evaluation gap a single-target benchmark cannot close.

\subsection{Summary and Remaining Gap}
\label{sec:mitig:summary}

\Cref{tab:mitigation} summarizes the per-category results (A100-SXM4, with GH200 efficiency); they separate into two kinds---the central result of RQ3.
The TileLang \emph{normalization and reduction family} are \emph{authoring artifacts}: the RC0 fix (\texttt{T.serial}$\to$\texttt{T.reduce} with native-dtype, tiled streaming I/O) brings every kernel to PyTorch-comparable or better (LayerNorm $124\%$, RMSNorm $990\%$, the rest $71$--$582\%$), closing gaps as large as $300\times$---on the A100 for the whole family, and on the GH200 for every family kernel but \texttt{logsumexp}, whose optimized config register-spills on \texttt{sm\_90} (a fixed kernel still carries a per-target tuning that must be re-validated).
The \emph{genuine residuals} survive best-effort optimization---Conv2d ($42\%$), large square GEMM (Matmul $77\%$ vs cuBLAS), and index-returning Argmax ($27\%$)---because on the data-center A100 cuBLAS, cuDNN, and \texttt{torch.argmax} are themselves well-tuned; of the Conv2d gap only ${\approx}2$--$3\%$ is absent Winograd selection (RC4), the rest general cuDNN implicit-GEMM tuning the Triton kernel does not match.

\paragraph{Cross-architecture portability is itself an authoring concern}
A \emph{correct, hand-optimized} DSL kernel can pass on one GPU yet collapse on another, so single-GPU or single-shape benchmarking cannot certify it.
\texttt{logsumexp} is one direction (fast on the A100, register-spilling on the GH200); the in-tree TileLang Softmax is the other, and purely an authoring choice: it caches the entire row in shared memory, so on the A100 its per-block footprint ($2N$ bytes at fp16) erodes occupancy---$20\times$ slower than PyTorch at $N{=}8192$ and $111\times$ at $N{=}32768$ (past the $48$~KB budget), while staying numerically correct---yet the identical kernel hits PyTorch parity on the GH200, whose larger shared-memory budget hides the defect.
The streaming variant (tiling the reduction over fixed $4$~KB blocks) holds parity across the sweep ($\rho{=}0.55$, \cref{tab:roofline}), so ``stream the reduction; do not cache the full row'' is a portable authoring pattern---and the collapse is a concrete instance of the evaluation gap, rooted in authoring rather than the language or hardware.

\begin{table}[t]
  \centering
  \caption{Summary of mitigation results. A100-SXM4 latencies in ms (lower is better);
    $E_\text{lib}$ = optimized-kernel efficiency relative to PyTorch/cuDNN on each GPU.
    \textbf{Bold} = ${\geq}95\%$ library efficiency on that GPU.}
  \label{tab:mitigation}
  \resizebox{\columnwidth}{!}{
  \begin{tabular}{lrrrrrl}
    \toprule
    & \multicolumn{3}{c}{A100-SXM4 (ms)} & \multicolumn{2}{c}{$E_\text{lib}$} & \\
    \cmidrule(lr){2-4}\cmidrule(lr){5-6}
    Kernel & PyTorch & Before & After & A100 & GH200 & Root cause \\
    \midrule
    \multicolumn{7}{l}{\textit{Triton}} \\
    \quad Matmul    & 0.607 & 1.081 & 0.788          & 77.0\%            & 68\%   & RC2 \\
    \quad Conv2d    & 3.44 & 17.86 & 8.12           & 42.4\%            & 34\%   & RC1+RC2 \\
    \midrule
    \multicolumn{7}{l}{\textit{TileLang}} \\
    \quad LayerNorm & 0.335 & 364  & \textbf{0.270} & \textbf{124\%}    & \textbf{120\%} & RC0 \\
    \quad RMSNorm       & 2.52  & 294   & \textbf{0.254} & 990\%$^\dagger$  & 1303\%$^\dagger$ & RC0 \\
    \quad Softmax       & 0.539 & 2.86  & 0.626          & 86.1\%           & 87\%   & RC0 \\
    \quad LogSoftmax    & 0.442 & 2.65  & 0.624          & 70.8\%           & 90\%   & RC0 \\
    \quad MaxReduction  & 0.560 & 11.97 & \textbf{0.418} & \textbf{131.6\%} & 84\%   & RC0 \\
    \quad MeanReduction & 0.768 & 10.65 & \textbf{0.766} & \textbf{99.5\%}  & \textbf{97\%} & RC0 \\
    % RESOLVED (2026-06-26 A100-SXM4 re-baseline): the A100 ncu sweep confirms logsumexp_opt does NOT spill on
    % sm_80 (0 local-mem load/store, 93 regs/thread, 30.5% occ; ncu_summary.csv) -> the 581.7% is genuine, an
    % sm_90-specific codegen divergence. See experiments/results/NVIDIA_A100-SXM4-40GB/NCU_FINDINGS.md.
    \quad LogSumExp     & 2.40  & 4.18  & \textbf{0.412} & \textbf{581.7\%} & 1.0\%$^\ddagger$ & RC0/RC3 \\
    \quad Argmax        & 0.622 & 11.97 & 2.306          & 27.0\%           & 22\%   & RC0+RC1 \\
    \bottomrule
  \end{tabular}
  }
  \vspace{0.3em}
  {\small $^\dagger$ RMSNorm $E_\text{lib}>100\%$: the fixed kernel fuses the scaling pass PyTorch's eager path leaves unfused.
  $^\ddagger$ LogSumExp's A100-tuned config register-spills on sm\_90 but not sm\_80; the $1.0\%$ (GH200) vs.\ $581.7\%$ (A100) split, retuned $63\%$ ceiling, and RC3 attribution are detailed in \cref{sec:mitig:other,tab:roofline}.
  \emph{Before} = unoptimized DSL kernel (Matmul: plain Triton at $4096^2$; Conv2d: baseline vs.\ implicit-GEMM Triton at $3\times3$); the norm/reduction \emph{Before} kernels are memory-latency-bound and clock-invariant. GH200 $E_\text{lib}$ is from \cref{tab:roofline} (unlocked).}

\end{table}
The \cref{tab:mitigation} mitigations are re-timed on the A100-SXM4 (locked clocks; the matmul figure is the $4096^2$ autotune value, reconciled against the $16384^2$ result in \cref{sec:mitig:other}), and the optimized kernels revalidate against the same per-dtype tolerances (5/5 pass, 0 edge-case crashes). RC2b-style search-space expansion is the higher-leverage of the two remaining convolution future-work directions.
